\section{zkInfer Design}
\label{sec:design}

\subsection{System Overview}
\label{sec:system-overview}

Building on Section~\ref{sec:characterization}, we design \TITLE{} as a
distributed runtime for serving zkML proofs under concurrent requests. A
request consists of a model and inference input. \TITLE{} applies the
configured partitioning policy, materializes the corresponding
per-partition inputs, and creates proof jobs that can execute across a pool
of prover nodes. The underlying proving protocol and proof-composition
mechanism remain unchanged.

Figure~\ref{fig:architecture} shows the architecture. The coordinator
maintains jobs from all active requests and assigns whole jobs to prover
nodes. Each node then controls when the assigned job's proving stages may
run. This separation is central to the design: the coordinator handles
request-level fairness, coarse placement, and setup locality, while the
node supervisor uses current local memory and stage-level estimates for
fine-grained admission.

\subsection{Partitioning and Self-Contained Proof Jobs}
\label{sec:self-contained-proving-jobs}

During request preparation, \TITLE{} partitions the model graph according
to a system-wide policy. Unless otherwise stated, we use one graph operator
per proof job. Section~\ref{sec:eval-sensitivity} evaluates coarser
partitions.

A partition may consume tensors produced elsewhere in the graph. \TITLE{}
therefore executes the model once before distributed proving to materialize
values that cross partition boundaries and attaches them to the consuming
jobs as witness inputs. Each job can then be scheduled independently of the
execution order of other jobs. Correctness does not rely on trusting these
values directly. Each partition is proved separately, while the backend's
existing proof-composition mechanism binds values produced by upstream
partitions to those consumed downstream~\cite{ezkl-splitting}.

\subsection{Cluster-Wide Job Placement}
\label{sec:global-placement}

The coordinator first chooses a request using round-robin fairness, then
selects a job from that request. Requests with no eligible job for the
polling node are skipped for that poll, and a request moves to the back of
the queue after an assignment if work remains. This prevents a large
request from monopolizing newly available prover capacity without requiring
preemption.

Let \(\widehat{m}^{\max}_j\) and \(\widehat{d}_j\) denote the predicted
maximum memory demand and duration of job \(j\). A job is ineligible for a
node if its known peak exceeds that node's physical memory. The coordinator
does not require the whole-job peak to fit within currently free memory,
because later memory-intensive stages may wait locally. Among eligible jobs
of the selected request, it first prefers jobs whose reusable setup
artifacts are already cached on the polling node, then selects the job with
the largest predicted duration, using submission order to break ties.

Stage-level packing remains local to each node. Keeping short-lived memory
conditions out of the coordinator avoids duplicating admission logic and
lets the node make decisions using the freshest view of its own usage.

\subsection{Node-Local Resource Admission}
\label{sec:local-admission}

Once a proof job is assigned to a node, the node supervisor controls memory
admission. Instead of reserving the job's predicted peak for its entire
execution, the supervisor admits each controllable stage according to that
stage's predicted demand and current local capacity.

\textbf{Stage-level estimates.}
For each model partition and proving configuration, the supervisor records
recent peak memory observations separately for each stage. With fewer than
five observations, \TITLE{} uses the largest observed value. Once at least
five are available, it uses the empirical 90th percentile of recent
history. If no stage-specific history exists, the controller falls back to
the job's whole-job estimate. Section~\ref{sec:profile-accuracy} evaluates
these predictions. Admission is memory-based because memory capacity
determines whether multiple proof stages can execute safely on the same
node.

Admission is enforced only at backend boundaries where execution can be
delayed without interrupting an operation already in progress. EZKL already
exposes major operations such as witness generation, setup, and proof
generation. We additionally instrument setup to distinguish parameter
loading, verification-key generation, and proving-key generation so that
resource demand can be measured at finer granularity.

\textbf{Reservation and live-memory checks.}
Let \(T_h\) be the physical memory of node \(h\), \(R_h\) a fixed system
reserve, and \(B_h=T_h-R_h\) the budget available to proving work. Before
stage \(s\) of job \(j\) begins, the supervisor obtains a predicted peak
\(\widehat{m}_{j,s}\). If \(A_h\) is the set of admitted stages, admission
requires
\begin{equation}
    \sum_{i\in A_h}\widehat{m}_i+\widehat{m}_{j,s}\leq B_h.
    \label{eq:stage-reservation}
\end{equation}

The supervisor also checks live available memory to account for memory
already resident in prover processes. Let \(r_i\) be the current resident
memory of an admitted process and \(F_h\) the operating system's available
memory after the reserve. The remaining protected growth is
\begin{equation}
    U_h=\sum_{i\in A_h}\max(0,\widehat{m}_i-r_i).
    \label{eq:unconsumed-reservations}
\end{equation}
For candidate stage \((j,s)\), let
\(g_{j,s}=\max(0,\widehat{m}_{j,s}-r_j)\). Admission additionally requires
\begin{equation}
    U_h+g_{j,s}\leq F_h.
    \label{eq:stage-live-memory}
\end{equation}
This avoids charging resident memory twice while preserving headroom for
expected growth.

Once admitted, a stage runs to completion without preemption. Its
reservation is released when the stage finishes, and the job must obtain a
new reservation before entering the next controlled stage. If memory is
insufficient, the job waits on the same node. A temporary startup
reservation covers the interval between leasing a job and its first normal
stage reservation.

\subsection{Reusable Setup Artifacts}
\label{sec:setup-reuse}

\TITLE{} persists reusable artifacts such as compiled circuits and proving
keys after their first generation. Each node maintains a bounded local
cache. A local hit avoids setup generation and restoration. If the artifact
exists only in shared storage, the node restores it locally. Otherwise it
is generated and persisted for later requests. During a cold execution,
\TITLE{} keeps newly generated proving material in memory across setup and
proving to avoid an unnecessary serialize-and-reload cycle.

Setup locality is a secondary placement objective. It is considered only
after request selection and resource feasibility, so a cache hit does not
override fairness or memory safety. When the cache exceeds its configured
limit, inactive artifacts are evicted in recency order. Artifacts used by
active jobs remain pinned until those jobs finish.

\subsection{Leases, Failures, and Recovery}
\label{sec:fault-tolerance}

Proof jobs are also the unit of recovery. Each execution attempt receives a
unique lease identifier. If a prover fails or a node stops responding, the
coordinator returns that job to the pending queue while leaving completed
jobs unchanged. Completion is accepted only when the reported lease matches
the job's current lease, so late results from an earlier attempt are ignored.
Heartbeats detect unavailable nodes, failed attempts are retried up to a
configured limit, and optional coordinator persistence restores completed
jobs while requeueing work that was running at restart.
